\section{Deployment, and the Gap It Exposed}
\label{sec:deployment}

\subsection{What happened}

The frozen stack was taken to the pool and did not act. In the first
runs the detector reported the anchor in a majority of frames, the
qualifier discarded every message, and the planner issued no lateral
command while the vehicle drove past the obstacle. Over one morning the
stack was changed until it avoided the anchor reliably. The pool was
available for that day only. Table~\ref{tab:gap} is the complete list of
what differed by the time the campaign ran, each entry verified against
both source files rather than recalled.

\begin{table}[t]
\centering
\caption{The deployment gap \realmark: every difference between the
stack that produced the frozen predictions and the stack that flew.
\emph{forced} marks a change made deliberately on the day;
\emph{unnoticed} marks a difference nobody acted on, which the
post-deployment ablation (Section~\ref{sec:diagnostic}) finds to matter
more than most of the forced ones.}
\label{tab:gap}
\setlength{\tabcolsep}{4pt}
\begin{tabular}{llll}
\toprule
quantity & frozen & flown & status \\
\midrule
qualifier warm-up updates & 20 & 2 & forced \\
qualifier confirm updates & 3 & 1 & forced \\
risk enter / exit threshold & 0.55 / 0.30 & 0.30 / 0.15 & forced \\
minimum avoidance hold & \SI{1.0}{\second} & \SI{4.0}{\second} & forced \\
engagement range & \SI{1.5}{\metre} & \SI{1.8}{\metre} & unintended \\
observation source & oracle relay + S1/S2 & image detector, blob
fallback & structural \\
approach distance & \SI{3.5}{\metre} & \SI{1.86}{\metre} & structural \\
observable range window & unbounded & \SIrange{0.29}{2.11}{\metre} &
unnoticed \\
clearing offset & \SI{2.5}{\metre} & \SI{2.5}{\metre} (basin
\SI{2.7}{\metre}) & unnoticed \\
run-past margin & \SI{4.0}{\metre} & \SI{4.0}{\metre} (basin
\SI{6}{\metre}) & unnoticed \\
run duration & \SIrange{90}{120}{\second} & \SI{25}{\second} & forced \\
\bottomrule
\end{tabular}
\end{table}

\subsection{The forced changes, and why each was necessary}

\emph{Qualification.} The warm-up requires twenty mutually coherent
updates. Simulated perception delivers a coherent detection every tick
and reaches twenty in a second; the real detector produced a box that
moved between frames, so twenty never accumulated and the qualifier
discarded all 1152 messages of a run while the vehicle drove past the
anchor. Two updates and one confirmation are what the real evidence
supports.

\emph{Risk threshold.} The risk score combines confidence, centrality and
apparent area. The simulated obstacle is a broad projection returned at
full confidence; the real anchor is a thin shank whose apparent area is
small, so the score stayed far below \num{0.55} even with the anchor
detected in \SI{79}{\percent} of frames and \SI{75}{\percent} of those
qualified. Lowering the threshold to \num{0.30} made the planner act;
\num{0.05}, tried first, made it react to the pool walls.

\emph{Avoidance hold.} A one-second hold was tuned where the simulated
vehicle strafes at \SIrange{0.20}{0.30}{\metre\per\second}. The real
vehicle reaches \SI{0.123}{\metre\per\second}, so one second buys
\SI{12}{\centi\metre} and the anchor is still there. Four seconds buys
about half a metre. This is the same shortfall that the S3 plant model
injects in simulation, appearing in the water.

\subsection{The differences nobody noticed, which mattered more}

\emph{The detector has a range window, and the experiment fitted inside
it.} The fallback detector accepts a candidate only if its image height
lies between \num{0.15} and \num{0.90} of the frame. Inverting
Equation~\ref{eq:range} at the \SI{90}{\degree} vertical field of view
the planner actually used, those gates are an observable range window of
\SIrange{0.29}{2.11}{\metre}. The vehicle started \SI{1.86}{\metre} from
the anchor, and the engagement range was \SI{1.80}{\metre}: six
centimetres apart. The physical planner was inside engagement range at
the instant of release. Nothing in the physical campaign measures
\emph{when} the planner decides to act, because the answer was fixed by
the geometry of the basin and the acceptance gates of the detector
before the vehicle moved.

That the approach was \SI{1.86}{\metre} was not a choice either. The
overhead camera used to measure the start pose covers
\SI{4.24}{\metre}$\times$\SI{2.38}{\metre}, so \SI{1.86}{\metre} is the
furthest the vehicle can start and still be wholly visible -- and a
start that cannot be measured before release makes a pre-registered
start region meaningless.

\emph{The manoeuvre does not fit the basin.} The committed planner
strafes to a \SI{2.5}{\metre} clearing offset and runs \SI{4}{\metre}
past the obstacle before returning. The basin is
\SIrange{2.7}{3.0}{\metre} wide, and the anchor is \SI{1.86}{\metre}
ahead of the start. Neither figure fits. Both were left at their frozen
values throughout the campaign, because the deployment effort was
directed at the parameters the calibration work had made salient. One
run ended against the far wall after a correctly executed avoidance.

\emph{The engagement range was never what the record said.} The session
record states \SI{5.0}{\metre}, with a justification for engaging on
sight. The runs used \SI{1.8}{\metre}: the wrapper that launched every
run forwards three arguments and the engagement range was the fourth, so
it silently fell back to the runner's default. All nine runs record
\SI{1.8}{\metre}. We report this because it is exactly the class of
error a matched sim-to-real study is supposed to catch and, in this
instance, only the final audit did.

\subsection{The planner that could not be flown}

The dynamic-window planner issued no command in the physical pipeline.
It evaluates route adherence and projects candidate velocities forward,
both of which need a pose estimate; in simulation this comes from
dead-reckoned odometry that the physical vehicle, with no velocity log
and no acoustic positioning, cannot produce.

The evidence deserves stating, because four of the six attempted
dynamic-window runs do contain a command stream and could be mistaken
for the planner working. Those streams take exactly three values in
surge, $\{0,\,0.08,\,0.15\}$, and three in sway,
$\{-0.2,\,0,\,+0.2\}$ -- which are the committed planner's nominal
surge, its reduced surge during avoidance, and its strafe setpoint. A
dynamic-window planner scores sampled velocity candidates and emits a
continuum; a three-level output is a finite-state controller's. Those
commands came from orphaned committed-planner processes surviving on the
same topic, the same contamination that the admission criterion of
Section~\ref{sec:realresults} exists to catch. The two runs with no
orphan present contain a command stream that is identically zero, and
that is the dynamic-window planner's actual output on this vehicle.

The consequence is not subtle and we do not soften it: half of the
frozen design was never tested, the physical campaign is
single-planner, and every comparison between the two planners in this
paper is simulated.
